\subsection{Strict changes under one blowup}\label{sec:nef-strict-blowups}

A blowup changes the amount of smooth volume near the vanishing set of a coefficient. A product equation turns that change into a strict increase of the degeneration order. Let $B=\mathbb P^m$, $m\geq2$, and normalize the Fubini--Study form $H$ by $\int_BH^m=1$. Let $\pi:\widetilde B\to B$ be the blowup of the affine origin. The incidence projection $\rho:\widetilde B\to\mathbb P^{m-1}$ gives $L=\rho^*\omega_{\rm FS}$; we also write $H$ for $\pi^*H$. The incidence embedding in $\mathbb P^m\times\mathbb P^{m-1}$ makes $\kappa=H+L$ K\"ahler.

\begin{proposition}\label{prop:nef-strict-blowup-orders}\index{Nef degeneration!strict blowup jump}
For every choice of integers $m\geq2$ and $p,k\geq1$, there is a single smooth-center blowup of a projective $(m+p)$-fold and a smooth semipositive nonclosed background whose common orders downstairs and after pullback are
\[
 a_{\rm down}=\max\{p-m/k,0\},\qquad a_{\rm up}=p-1/k.
\]
The orders exist on both sides. The increase can be any positive integer in varying dimension. The same family can lose positive limiting volume while both orders remain zero.
\end{proposition}

\begin{proof}
Put $W_0=\mathbb P^p$, with normalized Fubini--Study form $\omega$ satisfying $\int_{W_0}\omega^p=1$. Write $s=|z|^2$ on the affine chart of $B$. Choose a smooth $Q(s)\geq0$ equal to $s^k$ near zero, positive for $s>0$, and constant and positive for large $s$; it defines a smooth function $q$ on $B$. Put $n=m+p$, choose $b>0$, and define
\[
 \theta=bH+q\omega\quad\text{on }B\times W_0,\qquad
 \widetilde\theta=bH+(q\circ\pi)\omega
 \quad\text{on }\widetilde B\times W_0.
\]
The map $\Pi=\pi\times\mathrm{id}$ is the single blowup of the smooth center $\{0\}\times W_0$, and $\widetilde\theta=\Pi^*\theta$. Fix the K\"ahler directions $h=H+\omega$, $\widetilde h=\kappa+\omega$ and smooth positive probability densities $\mu=H^m\wedge\omega^p$, $\widetilde\mu=\kappa^m\wedge\omega^p/(2^m-1)$. Both backgrounds are smooth semipositive and nonclosed.

Write $\nu_{\rm down}=H^m$ and $\nu_{\rm up}=\kappa^m/(2^m-1)$ for the base probability forms, whose normalization is checked below. For either product write its perturbation as $\xi=A+g\omega$. Downstairs $A=(b+\epsilon)H$; upstairs $A=(b+\epsilon)H+\epsilon L$; in both cases $g=q+\epsilon$, with pullback understood upstairs. The form $A$ is K\"ahler on the relevant base. Use $\nu$ for the corresponding base form, and put $M=\int A^m$, $J=\int g^{-p}\nu$. The smooth equation \cref{thm:tw-smooth} and closed Stokes give $u$ with
\[
 \alpha=A+\ddc u>0,\qquad \alpha^m=(M/J)g^{-p}\nu.
\]
\Needspace{7\baselineskip}
The radial cumulative masses of the base forms are
\begin{equation}\label{eq:nef-radial-reference-masses}
 \begin{split}
 \Xi_{\rm down}(s)&=\left(\frac{s}{1+s}\right)^m,\\
 \Xi_{\rm up}(s)&=\frac{\left(1+\frac{s}{1+s}\right)^m-1}{2^m-1}.
 \end{split}
\end{equation}
To check these formulas and the constants $M$, write $t=\log|z|^2$ on the punctured affine chart. For a radial potential $F(t)$, direct polar integration of its Hessian gives
\[
 \int_{s_1<|z|^2<s_2}(\ddc F(\log|z|^2))^m
 =F'(\log s_2)^m-F'(\log s_1)^m.
\]
The potentials of $H$ and $L$ are $\log(1+\mathrm{e}^t)$ and $t$. Thus $xH+yL$ has endpoint slopes $y$ and $x+y$, and its total mass is $(x+y)^m-y^m$. The cases $(x,y)=(1,0),(1,1)$ give \eqref{eq:nef-radial-reference-masses}; the exceptional divisor and hyperplane at infinity have zero smooth volume. The form
\[
 \Omega=\alpha+g\omega=\xi+\ddc u>0,\qquad
\Omega^n=\binom{n}{m}\frac{M}{J}\nu\wedge\omega^p
\]
has constant volume density. The lower bound for all tests in \cref{thm:gaud-volume-dual-formula} is attained by
\[
 \gamma_*=g^{-p/(n-1)}\Omega.
\]
Indeed its $(n-1)$st power is
\begin{equation}\label{eq:nef-balanced-product-power}
 \gamma_*^{n-1}
 =\binom{n-1}{m}g^{-1}\alpha^m\wedge\omega^{p-1}
 +\binom{n-1}{m-1}\alpha^{m-1}\wedge\omega^p.
\end{equation}
The first coefficient multiplies a top-degree form on the base, and every factor in the second term is closed; hence both terms are closed. Thus $\gamma_*$ is balanced and Gauduchon. \Cref{thm:gaud-volume-dual-formula} gives the exact identities
\begin{equation}\label{eq:nef-strict-blowup-volumes}
 \mathcal V_\mu(\theta+\epsilon h)
 =\binom{n}{m}\frac{(b+\epsilon)^m}{J_{\rm down}},
 \qquad
 \mathcal V_{\widetilde\mu}(\widetilde\theta+\epsilon\widetilde h)
 =\binom{n}{m}
 \frac{(b+2\epsilon)^m-\epsilon^m}{J_{\rm up}},
\end{equation}
where $J_{\rm down}=\int_B(q+\epsilon)^{-p}\nu_{\rm down}$ and $J_{\rm up}=\int_{\widetilde B}(q\circ\pi+\epsilon)^{-p}\nu_{\rm up}$.

Near zero, differentiation of \eqref{eq:nef-radial-reference-masses} gives
\begin{equation}\label{eq:nef-radial-measure-asymptotics}
 \mathrm{d}\Xi_{\rm down}=(m s^{m-1}+O(s^m))\,\mathrm{d} s,\qquad
 \mathrm{d}\Xi_{\rm up}=\left(\frac{m}{2^m-1}+O(s)\right)\mathrm{d} s.
\end{equation}
\Needspace{10\baselineskip}
For a density $(a s^{d-1}+O(s^d))\,\mathrm{d} s$ and $q=s^k$, the singular integral has the following leading behavior as $\epsilon\downarrow0$:
\begin{equation}\label{eq:nef-radial-singular-integral}
 \begin{cases}
 \displaystyle\frac{a}{k}B(d/k,p-d/k)\epsilon^{d/k-p}
     &kp>d,\\[4pt]
 \displaystyle\frac{a}{k}\log(1/\epsilon)&kp=d,\\[4pt]
 \text{a finite positive limit}&kp<d.
 \end{cases}
\end{equation}
Here $B(u,v)=\int_0^\infty t^{u-1}(1+t)^{-u-v}\,\mathrm{d} t$. In the power case substitute $s=\epsilon^{1/k}v$, then $t=v^k$. At equality the integrand on $\epsilon^{1/k}\ll s\ll1$ is $a/s$. The $O(s^d)$ error is negligible: on a sufficiently small fixed interval its relative size is arbitrarily small, while the complementary integral is bounded. For $kp<d$, monotone convergence applies to the integrable limiting density. Since the numerators of \eqref{eq:nef-strict-blowup-volumes} tend to $\binom{n}{m}b^m>0$, formulas \eqref{eq:nef-radial-measure-asymptotics}--\eqref{eq:nef-radial-singular-integral}, with $d=m$ downstairs and $d=1$ upstairs, give the stated orders.

For $m=k=2$, $p=1$, the denominators satisfy $J_{\rm down}\sim\log(1/\epsilon)$ and $J_{\rm up}\sim(\pi/3)\epsilon^{-1/2}$; the volumes are respectively asymptotic to $3b^2/\log(1/\epsilon)$ and $(9b^2/\pi)\epsilon^{1/2}$. The orders change from $0$ to $1/2$. For $m=2$, $k=p=1$, the downstairs denominator has a finite positive limit, while $J_{\rm up}\sim(2/3)\log(1/\epsilon)$, so positive limiting volume downstairs becomes $9b^2/[2\log(1/\epsilon)]$ upstairs, with order zero on each side. For an arbitrary integer $j\geq1$, choose $m=j+2$, $p=j+1$, $k=1$. Then $p<m$ gives a positive downstairs limit, while the orders are $0$ and $j$.
\end{proof}
